\par\medskip
\begin{center}\textls[100]{b) Die allgemeine Theorie der elliptischen Funktionen.}\end{center}
\noindent\hspace*{\parindent}Über keine andere Periode von \textsc{Gauss'} wissenschaftlicher Laufbahn gibt
das \textit{Tagebuch} in so lückenloser Vollständigkeit Auskunft, wie über diese so
ertragreichen Frühlingswochen des Jahres 1800. Leider sind die zugehörigen
Aufzeichnungen nicht in derselben Vollständigkeit erhalten, so dass wir für
einige der im \textit{Tagebuch} angezeigten Entdeckungen genötigt sein werden, aus
späterer Zeit herrührende Ausarbeitungen heranzuziehen.

Wenn es \textit{Tagebuch} 105. vom 6. Mai 1800 heisst: »Die Theorie der transzendenten Grössen
\[
\int\frac{dx}{\sqrt{(1-\alpha x^2)(1-\beta x^2)}}
\]
\noindent haben wir zur vollsten Allgemeinheit durchgeführt«, so wird man hierunter
wohl die Aufstellung des allgemeinen Wertes zu verstehen haben, dessen
eine solche Transzendente für einen bestimmten Wert von \(x\) fähig ist, etwa
in dem Sinne, wie der allgemeine Wert von \(u=\int\frac{dx}{\sqrt{1-x^2}}\) durch \(u+2g\pi\),
der von \(u=\int\frac{dx}{\sqrt{1-x^4}}\) durch \(u+2(g+hi)\varpi\) für ganzzahlige \(g,h\) gegeben wird.
Gerade der Übergang von der für das lemniskatische Integral geltenden Form
\(u+2(g+hi)\varpi\) zu der Form \(u+2g\varpi+2hi\varpi'\) als dem allgemeinen Werte von
\begingroup\makeatletter\tagsleft@true\let\veqno\leqno\makeatother
\begin{equation*}
u=\int\frac{dx}{\sqrt{(1-x^2)(1+\mu^2x^2)}}\mathpunct{\vcenter{\hbox{,}}}\tag{17}
\end{equation*}
\endgroup
\noindent (dies ist die Normalform, die \textsc{Gauss} seinen hier zu besprechenden Untersuchungen zugrunde legt), mit
\begingroup\makeatletter\tagsleft@true\let\veqno\leqno\makeatother
\begin{equation*}
\frac{\pi}{M(1,\sqrt{1+\mu^2})}=\varpi,\quad
\frac{\pi}{\mu M\left(1,\sqrt{1+\frac1{\mu^2}}\right)}=\varpi',\tag{18}
\end{equation*}
\endgroup
\noindent stellt wohl die grösste Schwierigkeit dar, die \textsc{Gauss} bei den ihm zur Verfügung stehenden Hilfsmitteln zu überwinden hatte und wir wollen gleich
hier hervorheben, dass ihm eine ihn völlig befriedigende Einsicht in die Verhältnisse des Wertevorrats des Integrals einer \textls[100]{mehrdeutigen} Funktion einer\par
\clearpage\SourceLeaf{158}\ArithmeticHead{73}{R}{IVb. ALLGEMEINE ELLIPTISCHE FUNKTIONEN.}
\HistoricalFoot{X\textsubscript{2} Abh. 2.}{10}
\noindent komplexen Veränderlichen auch späterhin nicht vergönnt gewesen ist, und
dass wohl gerade dies einer der Gründe gewesen sein dürfte, die ihn von
einer Veröffentlichung seiner Ergebnisse abgehalten haben.

Für das lemniskatische Integral hatte er die Untersuchung in der Weise
geführt, dass er die Periodizität der Umkehrungsfunktion mit Hilfe des Additionstheorems bewies. Ähnlich verfährt er in einer Werke VIII, S. 93, 94
abgedruckten Notiz\footnote{Diese Notiz steht auf einem einzelnen Zettel, der sich in dem Päckchen mit Rechnungen befindet, auf die sich die Tagebuchnotiz 112. vom 12. Juni 1800 bezieht, und die Werke III, S. 426—431 abgedruckt sind. Der Zettel, der auch in Schrift und Tinte sich diesen Rechnungen anschliesst, wird also aus der hier in Rede stehenden Zeit stammen.} für die reelle Periode des Integrals \(\int\frac{dx}{\sqrt{1+x^3}}\mathpunct{\vcenter{\hbox{.}}}\) In der
Scheda Ac beginnen die Aufzeichnungen, die mit der Tagebuchnotiz 105. und
den darauf folgenden korrespondieren (siehe Werke X, 1, S. 194), mit dem
»\textls[100]{Problem}«, die durch die Differentialgleichung
\[
d\varphi=\frac{du}{\sqrt{1+\mu\cos u}}
\]
\noindent definierte Funktion \(\varphi\) von \(u\) und \(u\) von \(\varphi\) durch Reihen von der Form
\[
\begin{aligned}
\varphi&=u+\sum_{k=1}^{\infty}A_k\sin ku,\\[1ex]
u&=\varphi+\sum_{k=1}^{\infty}B_k\sin k\varphi
\end{aligned}
\]
\noindent darzustellen. Man könnte angesichts dieses nicht weiter ausgeführten Ansatzes an die Möglichkeit denken, dass \textsc{Gauss} zur Bestimmung des Wertevorrats des allgemeinen elliptischen Integrals erster Gattung einen ähnlichen
Weg eingeschlagen hat, wie \textsc{Jacobi} in seiner Abhandlung über die vierfachperiodischen Funktionen von 1834\footnote{\textsc{Jacobis} Werke II, S. 23; vergl. L. \textsc{Schlesinger}, \textit{Über Jacobis Auffassung des reellen Integrals als einer mehrdeutigen Funktion}, Bibliotheca Mathematica (3) XI, 1911, S. 138—152.}, es liegen aber keine sonstigen, diese
Vermutung bestätigenden Aufzeichnungen vor. Für das lemniskatische Integral \(u=\int\frac{dx}{\sqrt{1-x^4}}\) gibt \textsc{Gauss} ja im \textit{Tagebuch} 91a. und im \textit{Leiste} (Werke X, 1,
S. 149) die Entwicklung
\[
u=\frac\varpi\pi\varphi+\left(\frac\varpi\pi-\frac2\varpi\right)\sin2\varphi+\cdots,
\]
\clearpage\SourceLeaf{159}\ArithmeticHead{74}{L}{L. SCHLESINGER, ÜBER GAUSS’ FUNKTIONENTHEORETISCHE ARBEITEN.}
\noindent wo \(x=\sin\varphi\); dadurch wird in Evidenz gesetzt, dass \(u+2g\varpi\) für ganzzahliges
\(g\) zum Wertevorrat von \(u\) gehört, aber eine analoge Entwicklung, die \(u+2hi\varpi\)
ergäbe, findet sich nicht, und erst dies würde dem \textsc{Jacobischen} Gedanken entsprechen\footnote{Die auf die Umkehrung des Integrals
\[
\int\frac{dx}{\sqrt{(1-x^2)(1-\mu x^2)}}
\]
bezügliche Notiz Werke VIII, S. 96, die der Herausgeber \textsc{Fricke} (daselbst S. 97) mit unserer Tagebuchnotiz 105. in Verbindung setzt, steht auf dem Deckel eines Buches (\textit{Maupertuisiana}, Hamburg 1753), das \textsc{Gauss}, nach seinem eigenhändigen Eintrag, 1808 erworben hat, und kommt schon darum hier nicht in Betracht, weil die Periodizität der Umkehrfunktion darin überhaupt nicht erwähnt wird.}. \textsc{Gauss} dürfte wohl im allgemeinen Falle, ähnlich wie in den
oben genannten besondern, mit dem Additionstheorem operiert haben, was
auch durch die Stelle der Scheda Ac (Werke X, 1, S. 196, Art. 4) bestätigt
wird, wo für das elliptische Integral (17) das Additionstheorem aufgestellt
und zur Untersuchung des Verhaltens der Umkehrfunktion in bezug auf die
eine Periode \(\varpi\) benutzt wird. — Auch die oben erwähnte Bemerkung über
den von \textsc{Abel} eingeschlagenen Weg (Brief an \textsc{Bessel} vom 30. März 1828)
spricht für diese Auffassung, denn \textsc{Abel} hat in den \textit{Recherches}\footnote{\textsc{Abel}, Oeuvres ed. \textsc{Sylow-Lie}, I, S. 271.}, wie übrigens
auch \textsc{Jacobi} in den \textit{Fundamenta}\footnote{\textsc{Jacobis} Werke I, S. 85, § 19.}, die doppelte Periodizität mit Hilfe des
Additionstheorems bewiesen.

Auf die wirkliche Herstellung der Umkehrungsfunktion für das Integral
(17) bezieht sich die Tagebuchnotiz 106., wo es in direkter Anknüpfung an
die Notiz 105. heisst: »Am 22. Mai in Braunschweig gelang es, einen hochbedeutsamen Zuwachs dieser Theorie zu entdecken, wodurch auch alle früheren
Ergebnisse\footnote{Das bezieht sich offenbar auf die Theorie der lemniskatischen Funktionen! \textsc{Schlesinger}.}, ebenso wie die Theorie der arithmetisch-geometrischen Mittel
aufs schönste verknüpft und grenzenlos bereichert werden«. — Die fast überschwängliche Ausdrucksweise dieser Aufzeichnung lässt auf den tiefen Eindruck
schliessen, den die neu gewonnenen Einsichten auf \textsc{Gauss} ausgeübt haben.
Das »wahrhaft neue Feld der Analysis«, das er in den Tagebuchnotizen 92.,
95., 98. angekündigt hatte, war nun erschlossen und es übertraf alle gehegten
Erwartungen.

Wir finden einen Teil dieser Untersuchungen in der Scheda Ac von S. 26\par
\clearpage\SourceLeaf{160}\ArithmeticHead{75}{R}{IVb. ALLGEMEINE ELLIPTISCHE FUNKTIONEN.}
\HistoricalFoot{}{10\textsuperscript{\(\ast\)}}
\noindent an (Werke X, 1, S. 194ff.). Um allen Schwierigkeiten, die mit der direkten Betrachtung des Integrals (17) verbunden sind, aus dem Wege zu gehen, setzt \textsc{Gauss}
mit den durch die Gleichungen (18) definierten Grössen \(\varpi,\varpi'\), die sich für \(\mu=1\)
auf das lemniskatische \(\varpi=\frac\pi{M(1,\sqrt2)}\) reduzieren, den Ausdruck an
\begingroup\makeatletter\tagsleft@true\let\veqno\leqno\makeatother
\begin{equation*}
\operatorname{Sl}\psi\varpi=\frac\pi{\mu\varpi}\left(
\frac{4\sin\psi\pi}{e^{\frac12\frac{\varpi'}\varpi\pi}+e^{-\frac12\frac{\varpi'}\varpi\pi}}
-\frac{4\sin3\psi\pi}{e^{\frac32\frac{\varpi'}\varpi\pi}+e^{-\frac32\frac{\varpi'}\varpi\pi}}
+\text{etc.}\right),\tag{19}
\end{equation*}
\endgroup
\noindent der ganz nach der Analogie der für den Sinus lemniscaticus gefundenen
Reihe (1) Abschnitt III (S. 38) gebildet ist, und den \textsc{Gauss} in der nachher
zu besprechenden Tagebuchnotiz 108. den \textls[100]{Sinus lemniscaticus in allgemeinster Fassung} (universalissime acceptus) nennt\footnote{\textsc{Gauss} benutzt auch z. B. Werke X, 1, S. 200 (S. 39 der Scheda) das Zeichen \(\operatorname{sl}\) für diese allgemeine Funktion. Wir schreiben im Anschluss hieran, statt des von \textsc{Gauss} gewählten \(S\psi\varpi\), auch in der Gleichung (19) \(\operatorname{Sl}\psi\varpi\).}. Er setzt dann für
\[
\varphi=\psi\varpi,\quad\mu=\operatorname{tg}v,
\]
\begingroup\makeatletter\tagsleft@true\let\veqno\leqno\makeatother
\begin{equation*}
\operatorname{Sl}\psi\varpi=\operatorname{Sl}\varphi=\frac{T\varphi}{W\varphi}\mathpunct{\vcenter{\hbox{,}}}\tag{20}
\end{equation*}
\begin{equation*}
\left\{\begin{aligned}
T\varphi&=\sqrt{\frac\pi\varpi}\sqrt[4]{\frac1{\mu^2(1+\mu^2)}}
\left\{\frac{2\sin\psi\pi}{e^{\frac14\frac{\varpi'}\varpi\pi}}-\frac{2\sin3\psi\pi}{e^{\frac94\frac{\varpi'}\varpi\pi}}+\text{etc.}\right\}\\[1.5ex]
W\varphi&=\sqrt{\frac\pi\varpi}\sqrt[4]{\frac1{1+\mu^2}}
\left\{1+\frac{2\cos2\psi\pi}{e^{\frac{\varpi'}\varpi\pi}}+\frac{2\cos4\psi\pi}{e^{4\frac{\varpi'}\varpi\pi}}+\text{etc.}\right\},
\end{aligned}\right.\tag{21}
\end{equation*}
\endgroup
\[
\sqrt[4]{\frac1{1+\mu^2}}\sqrt{\frac\pi\varpi}=\sqrt{M(1,\cos v)},
\]
\noindent wo auch die Funktionen \(T\varphi,W\varphi\) nach der Analogie der für Zähler und
Nenner \(P,Q\) des Sinus lemniscaticus geltenden Reihenentwicklungen (3) Abschnitt III (S. 39) gebildet sind\footnote{Eine Verifikation des gleichzeitigen Bestehens der Gleichungen (19), (20), (21) gibt \textsc{Geppert}, a. a. O. \textsuperscript{83)}, § 38, S. 157.}.

Mit der Einführung der Funktionen \(T\varphi,W\varphi\) und der beiden weiteren\linebreak{}
\(T\left(\frac\varpi2-\varphi\right),\ W\left(\frac\varpi2-\varphi\right)\) hat \textsc{Gauss} die für die Theorie der elliptischen Funktionen genuinen Elemente geschaffen, dieselben, die \textsc{Jacobi} 1828 von sich aus
neu entdeckt und als \textls[100]{Thetafunktionen} bezeichnet hat. In spätern Aufzeichnungen (1825, Werke III, S. 465) benutzt \textsc{Gauss} neue Bezeichnungen,\par
\clearpage\SourceLeaf{161}\ArithmeticHead{76}{L}{L. SCHLESINGER, ÜBER GAUSS’ FUNKTIONENTHEORETISCHE ARBEITEN.}
\noindent die sich von denen, die in den gebräuchlichen Darstellungen der \textsc{Jacobischen}
Theorie angewandt werden\footnote{Vergl. z. B. H. \textsc{Weber}, \textit{Elliptische Funktionen und algebraische Zahlen}, Braunschweig 1891.}, nur unwesentlich unterscheiden. \textsc{Gauss} setzt:
\[
\begin{aligned}
x&=e^{-\pi\frac{\varpi'}\varpi},\quad y=e^{2\pi i\psi}=e^{2\frac\pi\varpi i\varphi},\\
P(x,y)&=1+x(y+y^{-1})+x^4(y^2+y^{-2})+x^9(y^3+y^{-3})+\text{etc.}\\
Q(x,y)&=1-x(y+y^{-1})+x^4(y^2+y^{-2})-x^9(y^3+y^{-3})+\text{etc.}=P(x,-y)\\
R(x,y)&=x^{\frac14}(y^{\frac12}+y^{-\frac12})+x^{\frac94}(y^{\frac32}+y^{-\frac32})+x^{\frac{25}4}(y^{\frac52}+y^{-\frac52})+\text{etc.}
\end{aligned}
\]
\noindent und zu diesen drei, in \(\psi\) geraden Funktionen tritt 1827 (Werke III, S. 472)
noch als vierte die ungerade Funktion
\[
S(x,y)=x^{\frac14}(y^{\frac12}-y^{-\frac12})-x^{\frac94}(y^{\frac32}-y^{-\frac32})+\text{etc.}=-iR(x,-y)
\]
\noindent hinzu. Ferner setzt er (Werke III, S. 465, 474)
\[
P(x,1)=p(x),\quad Q(x,1)=q(x),\quad R(x,1)=r(x)
\]
\noindent (siehe oben Gln. (\(\ast\)) S. 65) und nennt diese Reihen in einer Aufzeichnung in
seinem Handexemplar der \textit{Disquis. arithm.} (Werke III, S. 386) die »summatorischen Funktionen«. In den jetzt üblichen Bezeichnungen (\textsc{Gauss'} \(x\) entspricht dem \textsc{Jacobischen} \(q\)) ist
\[
\begin{gathered}
P(x,y)=\vartheta_{00}(\psi|x),\quad Q(x,y)=\vartheta_{01}(\psi|x),\\
R(x,y)=\vartheta_{10}(\psi|x),\quad S(x,y)=i\vartheta_{11}(\psi|x).
\end{gathered}
\]
